\begin{textcolorbox}[Prompt for Toy Case Judge]
\textbf{Task:} Evaluate the logical rigor of a solution to an inequality problem, focusing specifically on whether the direction of the inequality was justified using toy cases or special value substitution.\\

\textbf{Instructions:}

1. Carefully read the reasoning process used to solve the inequality.

2. Identify whether the direction of the inequality was determined by testing special values, trying toy cases, or relying on extreme-case analysis, rather than providing a general proof valid over the entire domain.

3. If the model uses a toy case (e.g., setting a variable to 0, 1, or choosing symmetric/equal values) or considers a variable tending to 0 or $\infty$ (extreme-case reasoning) to conclude the inequality direction, this should be flagged as logically unsound unless it is later supported by a rigorous or general argument.

4. Substituting special values for the purpose of verifying equality or testing sharpness is acceptable and should not be flagged.

5. If a toy case is used to show that the inequality does not hold (i.e., the two sides are incomparable), this is acceptable and should not be flagged.

6. Trying toy cases or substituting special values for the purpose of exploring or analyzing the problem—without using them to directly conclude the inequality direction—is acceptable and should not be flagged.

7. The goal is to confirm whether the final conclusion is justified for \textit{all variables in the given domain} by using sound and formal reasoning.\\

\textbf{Output Format:}

\texttt{<Analysis>:} Brief explanation of whether toy cases, special values, or extreme-case reasoning were used to justify inequality direction, and whether this undermines the rigor of the proof.\\
\texttt{<Flagged Reasoning Step (if applicable)>:} Quote or summarize the step(s) where toy cases or extreme cases were used to justify the inequality.\\
\texttt{<Answer>:} \texttt{True} or \texttt{False}. It indicates whether the solution's reasoning is logically valid (\texttt{True}) or contains flaws due to unjustified use of toy cases, special value substitution, or extreme-case reasoning (\texttt{False}).\\

\textbf{Key Considerations:}

1. Check which steps are used to justify the final general conclusion.

2. Toy case examples should not be used to conclude general inequality directions.

3. These techniques are allowed for illustration or equality testing, but not as a replacement for a full justification.

4. Extreme-case analysis used as a substitute for proof to justify the final general conclusion should be flagged.

5. Do not provide improvement suggestions---simply judge whether the logic used is valid or flawed.

6. Only reasoning steps that fail to generalize and are used to justify the final general conclusion should be flagged.\\

\textbf{Examples of Inputs and Outputs:}

\texttt{\{examples\}}\\

\textbf{Now analyze the following problem and solution:}

Original Problem:
\texttt{\{query\}}

Solution:
\texttt{\{response\}}

\end{textcolorbox}


